\subsection*{Universal Register Machine (URM)}

\mkImg{image6}

\hSep

\kwub{universal register machine (URM)} carries out the following computation,
starting with,

\begin{enumerate}
  \vItem Starting with \iMbox{R_{0}=0}, \iMbox{R_{1}=e} (\stu{code of a program}),

  \iMbox{R_{2}=a} (\stu{code of a list of arguments}) and \stu{all other registers zeroed}

  \vItem Decode \iMbox{e} as a RM program \iMbox{P}
  \vItem Decode \iMbox{a} as a list of register values \iMbox{a_{1}, \dots, a_{n}}
  \vItem Carry out the computation of the RM program \iMbox{P} starting with \iMbox{R_{0}=0, R_{1}=a_{1}, \dots, R_{n}=a_{n}}
  (and any other registers occurring in \iMbox{P} set to \iMbox{0})
\end{enumerate}


\hSep

\stub{Mnemonics for URMs} =>
\begin{enumerate}
  \vItem \iMbox{R_{0}} result of the simulated RM computation (if any)
  \vItem \iMbox{R_{1} \equiv P} \stub[orange]Program code of the RM to be simulated
  \vItem \iMbox{R_{2} \equiv A} list of RM \stub[orange]Arguments (or register contents) of the simulated machine
  \vItem \iMbox{R_{3} \equiv PC} \stub[orange]Program \stub[orange]Counter-label number of the current instruction
  \vItem \iMbox{R_{4} \equiv N} label number(s) of the \stub[orange]Next instruction(s)-also used to hold code of current instruction
  \vItem \iMbox{R_{5} \equiv C} code of the \stub[orange]Current instruction body
  \vItem \iMbox{R_{6} \equiv R} value of the \stub[orange]Register to be used by current instruction
  \vItem \iMbox{R_{7} \equiv S} and \iMbox{R_{8} \equiv T} are auxiliary registers.
  \vItem \iMbox{R_{9} \dots} other scratch registers.
\end{enumerate}

\hSep

\stub{Overall structure of the URM}
\mkImg{image5}

\hSep

